\documentclass[11pt,a4paper]{article}
\usepackage[T1]{fontenc}
\usepackage[utf8]{inputenc}
\usepackage{lmodern}
\usepackage[margin=25mm]{geometry}
\usepackage{microtype}
\usepackage[hidelinks]{hyperref}
\hypersetup{
  pdftitle={Re-reading some early works on free algebras for Lawvere theories},
  pdfauthor={Walter Tholen},
  pdfsubject={Online Algebra, Logic and Topology Seminar}
}
\setlength{\parindent}{0pt}

\begin{document}
\begin{center}
{\Large\bfseries Online Algebra, Logic and Topology Seminar}\\[4pt]
{\large (CMUC, University of Coimbra)}
\end{center}

\vspace{1em}
\textbf{Speaker:} Walter Tholen (York University, Canada)\\
\textbf{Date:} December 1st, 2026\\
\textbf{Time:} 16:00 (Europe/Lisbon)\\
\textbf{Place:} Google Meet

\vspace{1.5em}
\begin{center}
{\Large\bfseries Re-reading some early works on free algebras\\[3pt]
for Lawvere theories}
\end{center}
\vspace{1em}

\textbf{Abstract.} Our focus will be on Hugo Volger's construction of free algebras with respect to an algebraic theory, as given in his rarely cited 1968 paper [2] which, most likely, is based on his 1967 Diplomarbeit [1]. His construction is built on five general categorical lemmata of independent interest. Beyond Lawvere's thesis, these represent an early study of the preservation of finite products by left Kan extensions. In joint work with Mat\'ias Menni [7] we recently showed how Volger's construction extends to the left adjoints of algebraic functors. Other than this, in the talk we plan to present also some closely related results of the early papers [3,4] which directly refer to [2], as well as of the papers [5,6] which appear to have been written independently from [2].

\bigskip
\textbf{References}
\smallskip

\begin{list}{[\arabic{enumi}]}{%
  \usecounter{enumi}%
  \setlength{\leftmargin}{1.8em}%
  \setlength{\labelwidth}{1.4em}%
  \setlength{\labelsep}{.4em}%
  \setlength{\itemsep}{.4em}%
  \setlength{\parsep}{0pt}%
  \setlength{\topsep}{.3em}%
}
\small
\item H. Volger: \emph{Kategorien von Algebren \"uber algebraischen Theorien}. Diplomarbeit, Universit\"at Freiburg, 1967.
\item H. Volger: \emph{\"Uber die Existenz der freien Algebren}. Math. Z. 106:312-320, 1968.
\item D. Schumacher: \emph{Zur Existenz freier Algebren einer r-dimensionalen Theorie}. Manuscripta Math. 3:227-236, 1970.
\item C. Howlett and D. Schumacher: \emph{Free finitary algebras in a cocomplete cartesian closed category}. Canadian Math. Bull. 15(3):373-374, 1972.
\item F. Borceux: \emph{Universal algebra in a closed category}. Preprint, Univ. Catholique de Louvain, 1975.
\item F. Borceux and B. Day: \emph{On product-preserving Kan extensions}. Bull. Australian Math. Soc. 12:291-296, 1977.
\item M. Menni and W. Tholen: \emph{Revisiting Hugo Volger's paper \"Uber die Existenz der freien Algebren}. Preprint, 2026.
\end{list}
\end{document}
